%==============================================================================
% Capítulo 11: Recomendaciones y líneas de trabajo futuro
%==============================================================================
\section{Recomendaciones y líneas de trabajo futuro}\label{sec:recomendaciones}

A partir de los resultados, del análisis de errores (\autoref{subsec:analisis-errores}), de la calibración y la intervención humana medidas (\autoref{subsec:calibracion}, \autoref{subsec:reduccion-intervencion}) y de las limitaciones (\autoref{subsec:limitaciones-estudio}), se formulan recomendaciones para la evolución del sistema y para futuras replicaciones, en tres niveles de prioridad. Son prioritarias las que atacan los defectos medidos: la calibración del gate de confianza (\autoref{subsec:calibracion-gate}), el clasificador supervisado con corpus propio (\autoref{subsec:clasificador-supervisado}) y la corrección de la medición de intervención humana (\autoref{subsec:medicion-intervencion}). Fortalecen la evidencia o el clasificador el aprendizaje activo, la validación independiente, la replicación y los modelos abiertos. Amplían el alcance sin alterar la validez de lo medido las extensiones funcionales: canales, monitoreo, base de conocimientos, agente telefónico, la cola de respaldo del alta y las evoluciones de plataforma anticipadas en las Secciones~\ref{sec:arquitectura} y~\ref{sec:implementacion} (autenticación única corporativa, SSO; prefijo configurable del número de incidente; migración a Kubernetes).

\subsection{Calibración del gate de confianza}\label{subsec:calibracion-gate}

Por etapa (\autoref{subsec:calibracion}), el gate de 0,70 solo se evalúa sobre Gemini y la etapa de reserva, la etapa determinística se rige por su puntaje de correctitud (\autoref{subsec:subsistema-clasificacion}) y la confianza autodeclarada de Gemini, concentrada entre 0,85 y 1,00, no discrimina aciertos de errores, por lo que el gate no derivó ningún caso. Se propone, sobre un conjunto de validación propio, disjunto del corpus de evaluación: (i)~incorporar señales independientes del puntaje ---ambigüedad léxica entre categorías, margen entre las dos primeras categorías (que exige ampliar el contrato de salida, \autoref{anexo:H}), desacuerdo entre etapas y revisión por muestreo de casos de alta confianza (\autoref{subsec:human-in-the-loop})---; (ii)~recalibrar el puntaje por regresión isotónica, o por escalado de temperatura \parencite{Guo2017_calibration} si se obtienen puntajes por categoría, advirtiendo que una recalibración monótona corrige el nivel de la confianza verbalizada, sobreconfiada y dependiente de la elicitación \parencite{Xiong2024_confidence_elicitation, Tian2023_verbalized_confidence}, pero no crea capacidad de discriminación; y (iii)~derivar el umbral de una curva cobertura-precisión con ajuste fuera de la muestra, como el punto de operación de la etapa determinística, cuya ventaja frente a Gemini depende del vocabulario (\autoref{subsec:sensibilidad-ajustes}), con un piso declarado antes de observarla (la exactitud de la revisión humana o el 85~\% del segundo objetivo), y confirmarlo en un corpus independiente (\autoref{subsec:corpus-independiente}). La misma campaña permitiría contrastar los parámetros de inferencia, sin optimización documentada (\autoref{subsec:capa-procesamiento}).

\subsection{Incorporación de un clasificador supervisado con corpus propio}\label{subsec:clasificador-supervisado}

Se sugiere incorporar progresivamente un clasificador supervisado entrenado con el corpus etiquetado generado por el propio sistema durante su operación. La acumulación sostenida de decisiones validadas por los sectores responsables constituye un activo informacional valioso que, una vez alcanzado un volumen del orden de 5.000 casos (referencia indicativa de orden de magnitud, no derivada del corpus actual ni de un cálculo de potencia estadística), permitiría entrenar modelos de clasificación específicos del dominio organizacional. Las ventajas potenciales de esta migración incluyen: reducción del costo marginal por incidente clasificado al eliminar la dependencia del servicio externo de inferencia, disminución de la latencia al ejecutar la inferencia localmente, y ganancia de soberanía sobre los datos al evitar la transferencia internacional de fragmentos textuales. La disponibilidad de modelos de pesos abiertos, como Llama~2 \parencite{Touvron2023_llama2}, abre una vía hacia un sistema completamente autoalojado; sin embargo, en clasificación de tickets con muchas clases los modelos tipo BERT superaron a LLM abiertos como Llama y Mistral \parencite{paket2024itservice}, por lo que esa vía debería evaluarse con modelos ajustados al dominio y no solo mediante \emph{prompts}.

\subsection{Corrección de la medición de intervención humana}\label{subsec:medicion-intervencion}

La segunda sub-hipótesis quedó no contrastada por ausencia de registro primario (\autoref{subsec:reduccion-intervencion}, \autoref{subsec:hipotesis-trabajo}). Se propone registrar, para cada caso, la ocurrencia y el origen de la intervención (derivación por el gate, revisión por muestreo, corrección o rederivación por el sector receptor, completado manual de datos), una consigna de tiempo de inicio y de fin, el resultado (categoría confirmada o corregida) y el operador, seudonimizado. El registro ampliaría \texttt{clasificacion\_log}, que hoy guarda el indicador de revisión y la categoría validada, o se llevaría en una tabla dedicada vinculada al incidente, preferible cuando un incidente acumula varias intervenciones. La definición operacional seguiría las dos formas rotuladas de la \autoref{tab:variables} ---proporción de casos y proporción del tiempo del proceso--- y se aplicaría de igual modo en ambos brazos. Así podría contrastarse la sub-hipótesis contra la cota de 0,15, con el intervalo de Wilson de esa proporción, o como cobertura efectiva del gate, esto es, la proporción de clasificaciones erróneas que derivó a revisión. El mismo registro mediría las derivaciones erróneas y su costo en tiempo, insumos del balance neto (\autoref{subsec:balance-neto}).

\subsection{Cola de respaldo y alta asincrónica}\label{subsec:cola-respaldo}

La invocación de alta es sincrónica. Una evolución de mayor alcance la desacopla mediante un patrón \emph{outbox} o un intermediario de mensajes con reproceso, de modo que el alta se persista de forma diferida. Esa línea tensiona la respuesta sincrónica que el formulario web consume para informar el número de incidente y la medición de los tiempos de ingreso y persistencia (\autoref{subsec:desafios-integracion}) sobre la que se apoyan los resultados del Capítulo~\ref{sec:resultados}, y exigiría reintroducir la infraestructura de mensajería que el stack retiró (\autoref{sec:arquitectura}); se propone como trabajo futuro y no altera lo medido.

\subsection{Ampliación de canales de entrada}\label{subsec:ampliacion-canales}

Se propone ampliar el conjunto de canales de entrada para incluir mensajería corporativa instantánea, particularmente Slack y Microsoft Teams en aquellas organizaciones donde estas plataformas constituyen el canal de comunicación dominante. La inclusión de aplicaciones móviles propias para usuarios de campo constituye una extensión natural que atiende la demanda creciente de canales asincrónicos integrados al ecosistema laboral cotidiano. La arquitectura modular propuesta en la \autoref{sec:arquitectura} facilita esta ampliación: cada nuevo canal requiere únicamente un disparador adicional en N8N, sin modificaciones en las capas de procesamiento ni de persistencia.

\subsection{Panel de monitoreo en tiempo real}\label{subsec:panel-monitoreo}

El sistema incorpora un dashboard de analítica basado en ReCharts que permite visualizar tendencias históricas, distribución por sector y desglose por estado (véase \autoref{subsec:dashboard-analitica}). Como evolución de esta capacidad, se sugiere implementar una capa de monitoreo en tiempo real sustentada sobre tecnologías de observabilidad maduras, tales como Prometheus para la captura de métricas y Grafana para la visualización. Esta capa permitiría supervisar la latencia del sistema, la tasa de aciertos del clasificador, la distribución de carga entre sectores y el costo agregado de las invocaciones al modelo externo, habilitando alertas tempranas frente a degradaciones del servicio. La existencia de endpoints de salud (\texttt{GET /api/v1/health}) en cada componente del sistema proporciona los puntos de datos necesarios para instrumentar esta capa de observabilidad.

Más allá del monitoreo operativo, los datos acumulados de incidentes constituyen una fuente de inteligencia organizacional: su registro sistemático habilita flujos ETL (\emph{Extract, Transform, Load}) que permiten identificar patrones recurrentes de fallas, detectar áreas con mayor demanda de soporte y generar indicadores trimestrales que orienten decisiones estratégicas de inversión.

\subsection{Mecanismos de aprendizaje activo}\label{subsec:aprendizaje-activo}

Se propone incorporar mecanismos de aprendizaje activo (\emph{active learning}; \autoref{subsec:human-in-the-loop}) que seleccionen los casos para revisión y reentrenamiento mediante las señales independientes del puntaje de la \autoref{subsec:calibracion-gate}: ambigüedad léxica, margen entre categorías, desacuerdo entre etapas y muestreo de casos de alta confianza. La confianza declarada solo se usaría como criterio de prioridad una vez recalibrada y validada según esa subsección, porque en el rango observado de Gemini no discrimina aciertos de errores y la etapa determinística no emite una confianza utilizable. Las revisiones que alimenten el ajuste incremental serían las registradas según la \autoref{subsec:medicion-intervencion}, de modo que cada intervención mejore al clasificador con un costo operativo controlado.

\subsection{Validación sobre un corpus independiente}\label{subsec:corpus-independiente}

El clasificador evaluado se ajustó, en parte, con el mismo corpus sobre el que se reportan las métricas de clasificación (vocabulario del filtro determinístico y reglas de frontera del prompt), de modo que esas métricas son estimaciones dentro de la muestra. La confirmación del desempeño fuera de la muestra exige construir un corpus independiente, etiquetado con el mismo protocolo y sin solapamiento con el corpus de evaluación, y reportar sobre él las métricas de la etapa determinística y de la etapa semántica. Esa validación no se realizó en este trabajo porque no se dispuso de un corpus de preproducción con esas características.

\subsection{Estudios de replicación}\label{subsec:estudios-replicacion}

Se sugiere realizar estudios de replicación en organizaciones de distinto tamaño, sector y región geográfica, para fortalecer la validez externa de los hallazgos y reunir evidencia empírica sobre soluciones similares en contextos heterogéneos. Serían relevantes las réplicas en organizaciones del sector público (con más incidentes de sistemas administrativos), del sector industrial (con más incidentes de infraestructura física) y de mayor tamaño, donde el volumen diario supere las 100 unidades. Toda replicación debería medir además la exactitud del flujo manual sobre el mismo corpus, registrando por caso la derivación del operador, con etiquetadores distintos de los operadores del brazo manual (\autoref{subsec:limitaciones-estudio}), y cronometrar las reasignaciones (\autoref{subsec:balance-neto}). Convendría, asimismo, una asignación muestral no proporcional que sobremuestree las clases minoritarias ---como Bases de Datos, con un solo caso--- y pondere las estimaciones globales por la distribución del universo (\autoref{subsec:poblacion-muestra-corpus}).

\subsection{Evaluación de modelos de lenguaje de código abierto}\label{subsec:evaluacion-modelos-abiertos}

Se propone evaluar la integración de modelos de lenguaje de código abierto desplegados localmente ---tales como Llama~3, Mistral o Qwen--- como alternativa al modelo externo Gemini~3.6 Flash. Esta evaluación permitiría cuantificar el costo en exactitud frente a la ganancia en soberanía de datos, la eliminación de la dependencia de un proveedor único y la reducción de costos operativos a largo plazo. La realización de un benchmark comparativo sobre el mismo corpus de 200 casos constituiría una contribución adicional significativa al estado del arte.

\subsection{Base de conocimientos para resolución automática}\label{subsec:base-conocimientos}

El sistema se concentra en la clasificación y derivación de incidentes hacia los sectores responsables, pero no aborda la resolución de los casos que admiten soluciones estandarizadas y repetibles. Se recomienda incorporar una base de conocimientos, un repositorio estructurado de soluciones verificadas que, tras la clasificación inicial, permita proponer los pasos de resolución cuando el incidente coincida con un patrón documentado y validado. Ese autoservicio absorbería los casos simples y repetitivos ---reinicios de sesión, restablecimiento de contraseñas o configuraciones básicas de correo electrónico--- y reduciría la carga de los sectores de soporte.

Su construcción puede articularse con el aprendizaje activo (\autoref{subsec:aprendizaje-activo}): cada caso resuelto por un operador y validado como correcto sería candidato a incorporarse, previa curación editorial. La viabilidad técnica se apoya en la capacidad de los modelos de lenguaje para generar respuestas estructuradas a partir de documentación técnica, línea que conecta con la evaluación de modelos de código abierto (\autoref{subsec:evaluacion-modelos-abiertos}).

\subsection{Agente conversacional telefónico en tiempo real}\label{subsec:agente-conversacional}

Se propone, como evolución de mayor alcance del canal telefónico, reemplazar la captura de toma única ---una grabación de hasta 45 segundos que se transcribe de forma diferida--- por un agente que dialogue en tiempo real. Sobre Twilio Media Streams, que habilita el envío bidireccional del audio mediante WebSocket, podría formular preguntas de aclaración, comunicar el número de incidente asignado y cerrar la llamada, lo que elevaría la calidad de la captura en línea con el tercer objetivo específico.

Su adopción estuvo condicionada por la madurez de las interfaces de audio bidireccional, entonces en vista previa (\emph{preview}) y sujetas a cambios, y exige infraestructura no contemplada: un punto de entrada WebSocket con \emph{upgrade} en el proxy inverso, la transcodificación del audio (mu-law, 8~kHz) al formato del modelo (PCM, 16~kHz), la gestión de turnos de habla y un modelo de costo por minuto de transmisión, no por grabación. Se recomienda una prueba de concepto acotada que conserve la grabación y transcripción actuales como modo degradado de respaldo.

\newpage
